\documentclass{article}
\title{Pratice Problem 2.11}
\author{CSAPP}

\usepackage{booktabs}
\usepackage{listings}
\usepackage{parskip}
\usepackage{solarized-light}

\begin{document}
    \maketitle

        \begin{flushleft}
        \texttt{Armed} with the function inplace\_swap from
        Problem 2.10. you decide to write code that will reverse the elements
        of an array by swapping the elements from opposite ends of the array,
        working toward the middle.\\
        You arrived at the following function
            \smallskip
            \begin{lstlisting}{language=C}
                void reverse_array(int a[], int cnt){
                    int first, last;
                    for(first = 0, last = cnt - 1;
                    first <= last; 
                    first++, last--)
                    inplace_swap(&a[first], &a[last]);
                }
            \end{lstlisting}

            \section*{\small Solution A}
            For an array of odd lenght cnt = 2k + 1, what are the values of 
            the variables first and last in the final iteration of the function
            reverse\_aray?\\
            Answer\\
            Both first and last will have value of k, so we can attempt
            to swap the middle element.

            \section*{\small Solution B}
            Why does this call to function inplace\_swap set the array
            elements to 0.\\
            Answer\\
            This is because our function assummes x and y denotes different locations

            \section*{\small Solution C}
            What simple modification to the code for reverse\_array would 
            eliminate this Problem?\\
            Answer\\
            We have to replace the code in line 4 of reverse\_array to be 
            first < last. 
        \end{flushleft}

\end{document}